\documentclass[letterpaper,10pt,conference]{ieeeconf}

\IEEEoverridecommandlockouts
\overrideIEEEmargins

\usepackage{algorithm}
\usepackage{algpseudocode}
\usepackage{amsmath}
\usepackage{amssymb}
\usepackage{balance}
\usepackage{bm}
\usepackage{graphicx}

\title{\LARGE \bf
SCRIBE: Consistent Distributed Environment Modeling\\
for Communication-Limited Robotic Exploration
}

\author{Anonymous Authors}

\begin{document}

\maketitle
\thispagestyle{empty}
\pagestyle{empty}

\begin{abstract}
Autonomous exploration of the ocean and other communication-limited
environments requires agents to construct models from sparse measurements
while using those models to determine where to sample next. Distributed teams
can extend the scale and responsiveness of this process, but their ability to
learn collectively depends on exchanging information across intermittent,
bandwidth-constrained networks. Observation sharing grows with mission
duration and sensor dimensionality, while conventional model-state fusion can
become overconfident when model correlations are unknown. We introduce
SCRIBE, a fixed-dimensional distributed environment modeling process for
consistent model assimilation whenever communication is available. SCRIBE
represents a scalar field through Gaussian basis functions and exploits the
resulting linear information form to combine Information Kalman filtering,
Covariance Intersection, and Metropolis--Hastings consensus. Agents can
thereby exchange learned model states over any available peer link, without
knowledge of the full communication graph or waiting for scheduled
model-update windows. We evaluate SCRIBE with independent Vulcan
information-based planners in twelve simulated four-robot trials under
distance-limited communication and a complete blackout. The resulting
trajectories achieve pooled posterior variance within 3.2\% of ideal
centralized sharing, indicating that consistent model assimilation can
preserve effective distributed exploration through communication loss.
\end{abstract}

\begin{keywords}
autonomous exploration, distributed estimation, informative path planning,
intermittent communication, multi-robot systems
\end{keywords}

\section{Introduction}

Autonomous exploration is both a learning and a decision problem. A robot
entering an unfamiliar region must gather observations, construct a model of
the environment, and use that model to determine which future observations
will be most valuable. Model confidence validates what has been learned and
identifies where understanding remains incomplete. Gaussian Processes and related
uncertainty-aware models have consequently found broad use in adaptive
sampling \cite{hwang2019auv} and information-based path planning
\cite{cui2015mutual}.

\begin{figure}[t]
    \centering
    \includegraphics[
        width=\columnwidth,
        trim=0 0 0 0,
        clip
    ]{res/experiment_4/full/field_reconstruction.png}
    \caption{Simulated scalar field over a $10\,\mathrm{km}\times
    10\,\mathrm{km}$ domain (left) and one robot's final SCRIBE posterior
    mean (right). Large circles mark initial locations; dash--dot paths
    terminate in oriented robot symbols. The displayed model is the online
    posterior at the final sampling round, not a reconstruction formed by
    replaying the observation history.}
    \label{fig:field-reconstruction}
\end{figure}

The need for this capability is particularly clear in the deep ocean. An
autonomous underwater vehicle may be asked to localize a hydrothermal plume,
characterize a cold seep, or follow a chemical gradient using only the sparse
measurements gathered during its mission; autonomous cold-seep exploration is
one demonstrated example \cite{vrolijk2021using}. A team can observe a larger
region and pursue several scientific targets concurrently, but underwater
acoustic communication is intermittent, bandwidth-limited, and strongly
dependent on vehicle separation \cite{hwang2019auv}. Each robot must therefore
continue learning and planning without contact, then incorporate its peers'
learned model states whenever communication resumes.

Distributed estimates must remain probabilistically consistent. Existing work
shares gathered or synthetic observations
\cite{viseras2016decentralized,allamraju2017communication}, or distributes
Gaussian-process model updates \cite{jang2020multi}. Observation sharing grows
with the history and becomes especially burdensome for high-dimensional camera,
sonar, or spectrometer records. Learned states are compact, but states that
have interacted are correlated; treating them as independent repeatedly counts
the same observations and overstates certainty
\cite{chen2002estimation}. That failure directly affects exploration because
an overconfident model makes an unobserved region appear less informative than
it is.

Finite-dimensional approximations reduce the cost of distributed Gaussian
regression \cite{pillonetto2019distributed}. Related systems couple distributed
models to informative planning \cite{jang2021fully,nguyen2025spatially} or
target search \cite{kim2023distributed}. Together, these results motivate a
compact model-learning backend whose confidence remains consistent as
communication changes.

Signal-aware Collaboration for Reliable Information-based Estimation (SCRIBE)
makes three contributions. First, an Information Kalman Filter (IKF) makes
observations additive over Gaussian scalar-field (GSF) coefficients. Second,
Covariance Intersection (CI) reconciles divergent priors without their unknown
cross-correlations \cite{chen2002estimation,wang2012distributed}. Third,
Metropolis--Hastings Markov Chain (MHMC) weights distribute new information
using only the identities and degrees of presently available neighbors
\cite{xiao2005scheme}. Following the hybrid consensus--CI construction for
unreliable networks \cite{tamjidi2021unifying}, SCRIBE permits model fusion
whenever a peer link becomes available, without a full-graph view or a
scheduled team-wide update.

\section{Distributed Environment Modeling}

\subsection{Fixed-Dimensional Information State}

SCRIBE represents a scalar field through $n_\phi$ shared GSF bases
\cite{pillonetto2019distributed,jadaliha2012environmental},
\begin{align}
    \mathcal{F}(x,k)
    &=\bm{\psi}^{T}(x)\bm{\phi}(k),
    &
    \psi_i(x)
    &=\tau_i^{-1}
      e^{-\lVert x-\mu_i\rVert_2^2/\sigma_i^2}.
    \label{eq:field-model}
\end{align}
Only the coefficient vector $\bm{\phi}$ is estimated online. For
$\bm{z}_j=\bm{H}_j\bm{\phi}+\bm{v}_j$, agent $j$ maintains information matrix
$\bm{Y}_j=\bm{P}_j^{-1}$ and vector
$\bm{y}_j=\bm{Y}_j\hat{\bm{\phi}}_j$. A measurement contributes
\begin{align}
    \delta\bm{I}_j
    &=\bm{H}_j^T\bm{R}_j^{-1}\bm{H}_j,
    &
    \delta\bm{i}_j
    &=\bm{H}_j^T\bm{R}_j^{-1}\bm{z}_j.
    \label{eq:information-innovation}
\end{align}
The exchanged information vector and matrix are fixed by $n_\phi$, independent
of mission duration; replaying $k$ observations grows as $O(kd_z)$. Raw sensor
products can be processed locally before entering the shared field model.

\begin{algorithm}[H]
\caption{Distributed Scalar-Field Modeling at Agent $j$}
\label{alg:field-modeling}
\footnotesize
\begin{algorithmic}[1]
\Require GSF parameters $\{\mu_i,\sigma_i,\tau_i\}_{i=1}^{n_\phi}$
\State initialize $(\bm{y}_j,\bm{Y}_j)$ and
       $\hat{\mathcal{F}}_j(x,0)=\bm{\psi}^T(x)\hat{\bm{\phi}}_j(0)$
\While{executing the sampling mission}
    \State $(\bm{z}_j,\bm{X}_j)\gets\operatorname{GatherSamples}()$
    \State $(\bm{y}_j,\bm{Y}_j)\gets
           \operatorname{DistributedUpdate}(\bm{\psi},\bm{y}_j,\bm{Y}_j,
           \bm{z}_j,\bm{X}_j)$
    \State $\hat{\bm{\phi}}_j\gets\bm{Y}_j^{-1}\bm{y}_j$;
           $\hat{\mathcal{F}}_j(x)\gets
           \bm{\psi}^T(x)\hat{\bm{\phi}}_j$
\EndWhile
\end{algorithmic}
\end{algorithm}

\subsection{Fusion under an Unstable Network}

Current innovations are new evidence; priors can contain an unknown amount of
shared history. SCRIBE deliberately treats them differently. CI chooses
simplex weights that minimize the covariance of the fused neighbor priors
without assuming their cross-correlations. In parallel, innovations are
averaged with
$\gamma_{jm}=[1+\max(|\mathcal{N}_j|,|\mathcal{N}_m|)]^{-1}$.
These MHMC weights require only current neighbor degrees, rather than
knowledge of the global or future graph. After local convergence over a
reachable component of size $n_c$,
\begin{align}
    \bm{y}_j^{+}
    &=\widetilde{\bm{y}}_j+n_c\overline{\delta\bm{i}}_j,
    &
    \bm{Y}_j^{+}
    &=\widetilde{\bm{Y}}_j+n_c\overline{\delta\bm{I}}_j.
    \label{eq:hybrid-update}
\end{align}
CI prevents prior agreement from masquerading as a new observation; MHMC
propagates genuinely new information. This separation permits repeated,
peer-level fusion whenever links become available, rather than only during
scheduled team-wide update windows.

\begin{algorithm}[H]
\caption{SCRIBE Distributed Update at Agent $j$}
\label{alg:distributed-update}
\footnotesize
\begin{algorithmic}[1]
\Require Model $\bm{\psi}$; state $(\bm{y}_j,\bm{Y}_j)$; sample
         $(\bm{z}_j,\bm{X}_j)$
\State $(\bm{y}_j^{-},\bm{Y}_j^{-})\gets
       \operatorname{LocalPriorUpdate}(\bm{y}_j,\bm{Y}_j)$
\State $\bm{H}_j\gets\bm{\psi}(\bm{X}_j)$;
       form $(\delta\bm{i}_j,\delta\bm{I}_j)$ using
       \eqref{eq:information-innovation}
\State initialize prior and innovation consensus states
\While{local fusion has not converged}
    \State exchange model-state components and degree with neighbors
    \State fuse neighbor priors using CI
    \State average neighbor innovations using MHMC
\EndWhile
\State \Return $(\widetilde{\bm{y}}_j,\widetilde{\bm{Y}}_j)
       +n_c(\overline{\delta\bm{i}}_j,\overline{\delta\bm{I}}_j)$
\end{algorithmic}
\end{algorithm}

\noindent\textit{Computational scaling.}
Let $b$ be the GSF basis count or retained EOF rank, $N$ the agents, $E$ the
active links, and $d_j$ agent $j$'s degree. Each peer message is
$\Theta(b^2)$. MHMC costs $O(d_jb^2)$ per agent and $O(Eb^2)$ fleet-wide per
round; dense prediction costs $O(b^3)$, and $J_{\mathrm{CI}}$ CI evaluations
cost $O(J_{\mathrm{CI}}(d_jb^2+b^3))$ per agent. Thus bounded-degree fleets
scale linearly with $N$, whereas complete communication has $E=\Theta(N^2)$.
On a fixed connected graph, disagreement falls from $e_0$ to $\epsilon$ in
$\lceil\log(e_0/\epsilon)/(-\log\rho)\rceil$ MHMC rounds for mixing factor
$\rho<1$; under changing connectivity, its expectation depends on the link
process.

\begin{figure*}[t]
    \centering
    \includegraphics[
        width=0.97\textwidth,
        trim=0 0 0 0,
        clip
    ]{res/experiment_4/full/distributed_performance.png}
    \caption{Twelve simulated four-robot trials of 36 sampling rounds.
    Communication is distance- and bandwidth-limited to one in-range peer per
    robot in rounds 1--12 and 25--36; the shaded interval is a complete
    blackout. Curves show the mean and one standard deviation. The pooled-data
    panel isolates trajectory quality from fusion error. The disagreement axis
    is truncated to resolve the communicating methods; the annotation reports
    the no-communication peak. Lower values are better.}
    \label{fig:distributed-performance}
\end{figure*}

\section{Experimental Evidence}

We first evaluate SCRIBE independently of planning using four agents on fixed,
matched sampling paths. A continuous-sharing check reproduces the centralized
IKF to numerical precision. We then restrict every communication opportunity
to disjoint, distance-limited peer pairs, remove all links for the middle third
of the experiment, and restore only the same limited-range communication.
Across 30 simulated trials with a truth represented by the GSF model, SCRIBE
reaches final RMSE $0.1170\pm0.0114$, compared with $0.1007\pm0.0121$ for
centralized sharing, $0.1175\pm0.0115$ for CI alone, $0.2881\pm0.0204$ for the
IKF without CI, and $0.8136\pm0.0175$ without communication. CI prevents
priors formed during the blackout from being treated as independent when links
return, while the IKF innovation update recovers information that CI alone can
only conservatively intersect.

To demonstrate SCRIBE's effectiveness for exploration, four robots sample a
simulated field for 36 rounds under the same limited communication and blackout
schedule. Five comparisons isolate the sharing mechanism. IKF-only sums new
innovations without reconciling correlated priors; CI-only intersects local
posteriors without summing independent innovations; limited-observation
sharing sends the two newest unseen records over each available link. No
communication retains only local data, while centralized sharing supplies
every observation as an ideal bound. Every distributed method obeys the same
range limit: the largest component contains two robots, and the middle twelve
rounds contain no links.

We use Vulcan \cite{ayton2018vulcan} in a single-agent implementation for
simplicity. Each robot plans from its own current model, without joint
allocation or collaborative tree search. This deliberately isolates the
information reward supplied by SCRIBE to otherwise decoupled planners.
Figure~\ref{fig:field-reconstruction} shows the online reconstruction and
sampling paths. In Fig.~\ref{fig:distributed-performance}, SCRIBE's final
per-agent RMSE is $0.523\pm0.060$, compared with $0.458\pm0.052$ centrally,
$0.564\pm0.084$ for CI alone, $0.730\pm0.081$ for IKF alone,
$0.653\pm0.067$ for limited-observation sharing, and $1.041\pm0.053$ without
communication.

To evaluate exploration separately from fusion error, we pool the samples
selected under each regime and construct one common evaluation model.
Differences then arise from where the decoupled planners chose to sample, not
from how their online models fused information. The pooled RMSE is $0.455$ for
SCRIBE, $0.458$ centrally, and $0.594$ without communication. Integrated
posterior variance, the information reward used during planning, is $0.0488$
for SCRIBE, within 3.2\% of the centralized $0.0473$; CI alone reaches
$0.0528$, IKF alone $0.0617$, limited-observation sharing $0.0651$, and no
communication $0.1074$. SCRIBE therefore selects nearly as informative a set
of measurements as ideal sharing, even though no distributed planner ever has
access to a team-wide communication component.

SCRIBE communicates a learned state whose dimension is determined by $b$,
rather than by mission duration or raw-sensor dimension. Although the
simulated records are scalar, high-dimensional sonar, camera, or spectrometer
products can be processed locally before entering the learned field model.
Direct observation sharing must retain and transmit those records, a
distinction that grows throughout the mission.

\section{Conclusion}

SCRIBE provides a compact and probabilistically consistent environment model
for distributed exploration under limited communication. Fixed-path results
show that its hybrid update preserves consistency through disconnection and
reconnection; planning results show nearly centralized information collection
without a team-wide connected network. The communicated state remains
$O(b^2)$, bounded-degree fleet work scales linearly with the number of agents,
and local models remain usable throughout a blackout.

Because its update is independent of the planning architecture, SCRIBE can
serve fully distributed GP and MCTS systems \cite{jang2021fully} and
coordinated planners such as Multi-Agent Vulcan \cite{olkin2024multi}. The
same requirements arise in ocean, planetary, subterranean, and other remote
missions where neither teammates nor human operators are continuously
reachable.

\balance
\bibliographystyle{IEEEtran}
\bibliography{scribe_workshop_paper_references}

\end{document}
